%% Copyright 2012 Renater
%
% This work may be distributed and/or modified under the
% conditions of the LaTeX Project Public License, either version 1.3
% of this license or (at your option) any later version.
% The latest version of this license is in
%   http://www.latex-project.org/lppl.txt
% and version 1.3 or later is part of all distributions of LaTeX
% version 2005/12/01 or later.
%
% This work has the LPPL maintenance status `maintained'.
%
% The Current Maintainer of this work is Matthieu Herrb
%
% This work consists of the files:
% jres.cls
% article.tex
%----------------------------------------------------------------------
%
% Instructions pour utiliser la classe jres:
%
% Suivre cet exemple au maximum. Règles générales:
%
% - suivre les règles de saise de french/babel concernant les
%   ponctuations doubles
% - le style de bibliographie est imposé, pas besoin de le re-définir.
%
%----------------------------------------------------------------------
% Utilisation de la classe de documents jres.
% Les parametres city et year permettent de situer le lieu et l'année
% (optionnels si la version à jour de la classe est utilisée
\documentclass[city=Montpellier,year=2013]{jres}

    % si pdflatex est utilisé, spécifier le codage du fichier source
    \usepackage[utf8x]{inputenc}

    \title{Titre de l'article JRES 2013}

    % Utiliser la commande \auteur pour definir la liste des auteurs.
    % Le 3e élément (adresse) peut être vide

    \auteur{Prénom Nom de l'auteur 1}
	{Laboratoire ou service...}{%
	    1, rue de l'adresse\\
	    34\,000 Montpellier}

    \auteur{Prénom Nom de l'auteur 2}
	{Laboratoire ou service}{%
	    2, rue du second auteur\\
	    34\,000 Montpellier}

    % Pas besoin de \date{}.

\begin{document}

% produit le titre selon les infos ci-desses
\maketitle

% Résumé
\begin{abstract}
Ce document résume les instructions aux auteurs de la conférence JRES~2013.
Il est téléchargeable sur le site des JRES : \url{http://www.jres.org}.

Le résumé doit faire de quinze à vingt lignes. Il est écrit en italique,
utilsant une police avec empattement en taille 10 points.
Le titre « Résumé » (ainsi que le titre « Mots-clefs » ci-dessous)
sont en 14 points, gras, sans empattement.
\end{abstract}

\motscles{IPv6, PKI, CMS...}

\section{Introduction}

Voici les instructions de présentation pour la rédaction d'un article
pour les JRES en utilisant \LaTeX{}.

Il est recommandé d'utiliser un moteur \LaTeX{} produisant directement
du PDF (\texttt {pdflatex}, \texttt {xelatex} ou \texttt {luatex}) et il est
impératif d'utiliser le format « Portable Document Format » (.pdf) pour
soumettre un article au comité de programme.

Les auteurs doivent utiliser la classe de document « jres » téléchargeable
sur le site des JRES (\url{https://2013.jres.org/auteurs/modeles}).

\section{Présentation générale}

\subsection{Structuration et polices}

La police de base est de type \emph{Times} pour le corps du texte, et
\textsf{Helvetica} pour les titres.
Le texte  normal est en 10 points.

La police pour les extraits de fichiers de configuration ou de code
source est \texttt{Courier} en 10 points.

L'environnement pour les extraits de fichiers de configuration ou
de code source  est « lstlisting ». Voici un exemple :
\vspace{2ex}
\begin{lstlisting}
10.0.0.1   serveur-1
127.0.0.1  localhost.localdomain localhost
\end{lstlisting}

Les titres de section sont en 14 points gras sans empattement
(Helvetica ou Arial), avec numérotation arabe.

Les titres de sous-section sont en 12 points gras sans empattement,
avec numérotation arabe (avec rappel du numéro de section et
séparation par un point).

Les titres de sous-sous-section sont en 10 points gras sans empattement,
avec numérotation arabe (avec rappel du numéro de sous-section et
séparation par un point).

Les titres de niveau inférieurs sont à éviter, pour des raisons de
lisibilité.
Les titres sont en couleur bleue sombre (valeurs RGB : 0.,0.2,0.375).  

Les notes de bas de page sont en 8 points\footnote{Comme ceci}.

\subsection{Typographie}

Les auteurs s'attacheront, dans la mesure du possible, à respecter les
règles typographiques françaises.

En particulier, ils veilleront à utiliser des espaces insécables avant
les ponctuations hautes (;:!?) et des guillemets français (« »).

Il est suggéré aux auteurs de lire les leçons de
typographie \cite{andre1990} de Jacques André.

\subsection{Bibliographie}

La bibliographie sera faite avec des numéros et dans l'ordre
d'apparition dans le texte en utilisant le style de bibliographie
\emph{jres} (voir fichier « article.bib »).

\begin{enumerate}

    \item Pour les articles d'une revue, on indiquera obligatoirement les
      auteurs, le titre de l'article, la revue, le numéro, le volume,
      l'année et le mois. Voir \cite{exemple1}

    \item Pour les articles dans les actes d'un congrès, on mentionnera
      obligatoirement les auteurs, le titre de l'article, le titre des
      actes, l'année. Voir \cite{exemple2} (le mois et le lieu sont
      facultatifs mais conseillés)

    \item Pour les livres, on citera obligatoirement le ou les auteurs, le
      titre, l'éditeur et l'année. Voir \cite{exemple3} (il est conseillé de
      rajouter le lieu d'édition)

\end{enumerate}

\subsection{Adresses électroniques}

Pour diminuer les risques de moissonnage, aucune adresse mail ne
figurera dans l'article. 

Pour les références à des pages web (sauf s'il s'agit d'une référence
auquel cas elles doivent être placées dans la bibliographie), la macro
\verb|\url| sera utilisée, par exemple comme ceci : \url{http://2013.jres.org/}

\section{Mise en page}

\subsection{Nombre de pages}

Dans le cas des posters, la taille doit être comprise entre une et
sept pages. Le poster lui même sera inséré avec cet article ; pour les
démonstrations et les présentations courtes, l'article doit comprendre
entre deux et huit pages ; pour les présentations longues, la taille
de l'article doit être comprise entre six et huit pages.

\subsection{Marges}

Les marges sont les suivantes :

\begin{itemize}
    \item en haut : 2\ cm ;
    \item à gauche et à droite : 2\ cm ;
    \item en bas : 2\ cm.
	\begin{itemize}
	    \item Liste à puce de second niveau
	    \item 2e élément
	\end{itemize}
\end{itemize}

Les pages doivent être numérotées. Le pied de page du modèle doit être
maintenu.

\subsection{Figures}

Les figures sont de préférence insérées à proximité du texte qui y
fait référence. Elles sont numérotées et suivies d'une légende
explicite, en italique, non gras. Le texte fait référence à la
figure~\ref{fig-exemple}.

Le paquet \texttt{graphicx} est utilisé pour inclure des images.

\begin{figure}
    \centerline{\includegraphics[height=4cm]{figure}}
    \caption{Un exemple de figure}
    \label{fig-exemple}
\end{figure}

En cas d'insertion de figure contenant des couleurs, il est souhaitable
d'en assurer la lisibilité sur une imprimante noir et blanc.

\section*{Annexe}

Les auteurs ajoutent, si nécessaire, une annexe avec une section non numérotée.

% Bibiliographie
\nocite{*}
\bibliography{article}

\end{document}
